\section{Threat Model: SEU Emulation in Volatile Memory}\label{sec:threat}

The faithfulness theorem (\cref{sec:faithfulness}) is meaningful only
against a precisely stated perturbation family. This section
establishes that family. We separate three concerns: the canonical
SEU primitive that RadFI emulates (\cref{ssec:seu-primitive}), the
kernel transit path on which RadFI operates
(\cref{ssec:kernel-transit}), and the DRAM-bidirectional refinement
introduced in v0.1.2 following empirical observation of the Linux
7.0.x I/O path (\cref{ssec:dram-bidir}).

\subsection{The Canonical SEU Primitive}\label{ssec:seu-primitive}

Following the canonical SEE
taxonomy~\cite{dodd2003seu,baumann2005softerrors}, a Single-Event
Upset is a single-bit (occasionally multi-bit) state transition of a
memory cell induced by the charge deposited by an ionizing particle.
The probability of a SEU striking a given cell exposed to a particle
flux \(\Phi\) of linear energy transfer LET is, to first order,
\(P_{\text{SEU}} = 1 - \exp(-\sigma(\text{LET}) \, \Phi \, t)\),
where \(\sigma(\text{LET})\) is the cross-section curve of the cell
technology~\cite{petersen1996crosssection}.

For the purposes of RadFI, the modelled primitive is the
\emph{single bit-flip} --- a transition of one bit from 0 to 1 or 1
to 0, in a memory location reached by a kernel I/O path. We do not
model the underlying physics; we model the \emph{observable
consequence} at the byte level. This is the same abstraction used by
software-implemented fault injection tools in the canonical
literature~\cite{carreira1998xception,hsueh1997faultinjection}, and
is the level at which a filesystem recovery operator can defend.

\subsection{Kernel Transit Path}\label{ssec:kernel-transit}

The action surface of RadFI is the bio-layer of the Linux block I/O
path. A \texttt{struct bio} is the kernel's representation of an
in-flight I/O request; it carries a vector of \texttt{bio\_vec}
segments, each pointing to a kernel memory page with offset and
length. When a filesystem issues a read or write, the kernel
constructs one or more bios and submits them through
\texttt{submit\_bio\_noacct()} or, in some buffer-cache paths,
\texttt{submit\_bh()}~\cite{corbet2005ldd3}.

RadFI intercepts these submission entry points through
kprobes~\cite{corbet2017kprobes,kernel-kprobes}. At the moment of
interception, the bio's payload is in DRAM, in transit to or from
the storage device. A single bit-flip applied to one byte of the
first \texttt{bio\_vec} segment is the operational realization of an
SEU on this payload.

\subsection{The DRAM-Bidirectional Model}\label{ssec:dram-bidir}

RadFI v0.1.0 modelled bit-flips on \emph{write} bios only, on the
rationale that a flip on a read bio would be re-fetched from storage
and effectively self-correct. Empirical observation during the
BEAMFS v2 development on 2026-04-29 invalidated this rationale:
under Linux kernel 7.0.x, a read of an \texttt{sb\_bread}-backed
block (the BEAMFS metadata path) traverses
\texttt{\_\_bread\_gfp} \(\to\) \texttt{\_\_getblk\_slow} \(\to\)
\texttt{bh\_read} \(\to\) \texttt{submit\_bio\_noacct}, and the
resulting payload is delivered to the page cache without re-fetch. A
flip on a read bio is therefore \emph{persistent} until the next
page cache eviction, and corrupts the consumer of the data exactly
as a flip on a write bio corrupts the disk.

RadFI v0.1.2 introduces the \emph{DRAM-bidirectional model}
(intermediate v0.1.1 added internal counter accounting for F5 only):
bit-flips can target both write bios (modelling SEU on memory
transit toward storage) and read bios (modelling SEU on memory
transit from storage to user). The choice is governed by the
\texttt{inject\_on\_read} debugfs flag (default true). Setting it to
false restores the v0.1.0 write-only behaviour.

This model is more conservative --- and more physically realistic
--- than the v0.1.0 model: a DRAM cell holding bio payload is
vulnerable to SEU regardless of the data flow direction, and an
instrument that injects only one direction misses half the empirical
surface.
